\section{Introduction}

Chemical exchange between the flowing aqueous phase and the solid matrix
governs the transport of dissolved species in porous media alongside advection
and diffusion. In many subsurface systems this exchange proceeds by surface
complexation, ion adsorption, and fast local equilibrium reactions, which
retard dissolved species, reshape concentration fronts, and alter the apparent
accessibility of pore space \citep{dzombak1990surface}. The effects are
strongest in multi-scale rocks that combine open macropores with microporous
matrix regions, where a single image may contain fast transport pathways, slow
exchange domains, and a large internal reactive area that the image cannot
resolve \citep{menke2021upscaling}. The same physics governs contaminant
retention and acid mine drainage treatment
\citep{webster1998trace,blowes2000treatment}, critical-mineral recovery and
direct lithium extraction \citep{reich2022lithium,vera2023environmental}, and
radionuclide transport for geological disposal
\citep{miller2010upscaling,dewindt2019reactive,bildstein2019rtm}, in all of
which selective adsorption or ion exchange may be controlled by surface area
that is not resolved at the image scale.

Three classes of model are available, and each is limited in a different way.
Classical Darcy-scale reactive transport models represent adsorption and
surface complexation through effective source terms, bulk sorption
coefficients, or equilibrium closures applied to the continuum volume
\citep{steefel1990kinetic,steefel1994coupled,steefel2015reactive,mayer2002multicomponent,xu2006toughreact,lichtner1985continuum}.
They remain indispensable for large-domain prediction, but they require
constitutive upscaling assumptions for both transport and reactive surface
area, and they cannot reproduce the spatial heterogeneity of surface inventory
that image-based data make visible \citep{battiato2011applicability}. Fully
resolved pore-scale models, including direct numerical simulation
\citep{molins2021simulation,molins2015reactive,molins2012investigation,starchenko2018development},
lattice Boltzmann methods \citep{yoon2015lattice,kang2007lattice}, and
explicit-boundary DBS approaches
\citep{maes2018direct,maes2021geochemfoam,soulaine2017mineral}, instead locate
reactive uptake at the meshed grain surface. This is appropriate when the
dominant surface area coincides with the resolved interface between pore and
grain, but it is bounded by image resolution: once internal micro-scale area
falls below the mesh scale, a boundary-resolved description is no longer
faithful to the physical geometry, and the cost of resolving it rises steeply.
Pore-network and other reduced pore-scale models occupy a middle ground
\citep{blunt2001flow}, but they depend on idealised throat-and-pore
representations and do not deliver the voxel-by-voxel field description used in
image-based DBS simulations.

The consequences are specific in image-based simulations of carbonate and other
multi-scale porous media \citep{menke2015dynamic, menke2018fourd,
menke2021upscaling,wildenschild2013xray,blunt2013porescale}. Segmented images
routinely separate macropore, microporous, and nearly solid regions, but the
reactive internal surface of the microporous phase is not the same object as
the explicitly meshed boundary of the coarse computational cells. Ignoring that
internal area under-predicts reactive uptake. Inserting it as an ad hoc
volumetric sink leaves both the conservation properties and the meaning of the
cell-scale surface inventory unclear.

PHREEQC and its reaction-module form, PhreeqcRM, already provide the required
chemistry: aqueous speciation, surface complexation, equilibrium and kinetic
reactions, and operator-split coupling to transport simulators
\citep{parkhurst2013phreeqc, parkhurst2015phreeqcrm,charlton2011modules}.
Existing PHREEQC-coupled workflows serve structured-grid continuum models and
operator-split multi-component solvers well \citep{jyoti2021validation}. The
open question is how to embed that engine in a transport formulation consistent
with unresolved image-based geometry. The closest existing approach is
\texttt{porousMedia4Foam} \citep{soulaine2021porousmedia4foam,
pavuluri2022reactive}, which couples a micro-continuum OpenFOAM solver to
PHREEQC at pore, hybrid, and Darcy scales, but which applies chemistry through
kinetic or equilibrium flux conditions at the diffuse fluid-solid interface. It
therefore inherits the limitation of the fully resolved methods, in that the
reactive surface must overlap the resolved computational boundary, and it draws
no explicit distinction between the numerical DBS porosity used to regularize
transport and the physical mobile-water fraction that governs chemistry.
Earlier PHREEQC-coupled cases within GeoChemFoam itself
\citep{maes2022heat,maes2022vos} share this restriction: they support reactive
transport and surface complexation where the reactive surface coincides with an
explicit mesh boundary, which remains valid but does not cover the case in
which the dominant reactive area lies inside the coarse cell.

In this work we develop a DBS-compatible reactive transport formulation in
which the reactive surface is a property of the cell rather than of the mesh
boundary. Aqueous transport is written for the physical mobile-water fraction;
unresolved reactive surface area is prescribed as a cell field and carried as a
cell-centred adsorbed inventory; local PHREEQC chemistry transfers mass
conservatively between the aqueous and adsorbed states; and no heuristic
redistribution of source terms is used. Two consequences follow. The
formulation requires an explicit separation between the numerical DBS porosity
that stabilizes transport in nearly solid voxels \citep{soulaine2016micro} and
the physical mobile-water fraction that sets chemical concentrations. It also
requires the surface inventory to be defined at cell centres and driven by an
unresolved internal surface-area density, rather than reconstructed on mesh
faces. Together these choices let the method act on microporous matrix regions
instead of only on explicitly resolved boundaries between pore and grain. The
governing equations are posed on general finite-volume meshes; the image-based
cases used here are the motivating application.

The solver that results, \texttt{reactiveTransportDBSFoam}, couples
conservative DBS transport to PHREEQC-based surface complexation at the cell
scale. A one-dimensional analytical retardation solution provides a closed-form
reference for the coupled advection, diffusion and sorption system, and
closed-system batch calculations test the aqueous and adsorbed mass balance.
The surface-area algorithms are verified on controlled geometries of known area
and on image-based micromodels, and a resolved-to-DBS two-scale benchmark of
selective Li$^{+}$ uptake tests whether the upscaled cell-based description
reproduces a fully resolved calculation of the same system. The formulation is
then applied to three image-based geometries of increasing complexity:
Pb$^{2+}$ uptake on hydrous ferric oxide (HFO) in a two-dimensional fracture
network, as an analogue for heavy-metal remediation in fractured aquifers and
permeable reactive barriers; selective Mn$^{2+}$ uptake from a Mg$^{2+}$-rich
brine on HFO in a three-dimensional Kelvin-cell sorbent, as an analogue for
critical-metal capture; and Sr$^{2+}$ migration in Estaillades limestone, as an
analogue for radionuclide transport in a deep geological repository (DGR).
Together these cases determine how resolved pore geometry and unresolved
internal surface area interact to control the spatial distribution of aqueous
and adsorbed species in multi-scale porous media.
